% =======================================================
% 6. RESULTS
% =======================================================
\section{Results and Detailed Analysis}
\label{sec:results}

\subsection{Pipeline reduction on application-scale code}
\label{sec:rq-field}
Table~\ref{tab:pipeline} follows the loops of the 200-file corpus through the deterministic phases. Phase~1 finds \plRaw{} natural loops, of which \plHeader{} lie in system and library headers (mostly STL containers) and are discarded. Of the remaining \fieldLoops{} project loops, HF-Static decides \plDecided{} (\plDecidedPct{}) without any model: \fieldPar{} are proven parallel and \fieldSer{} serial, the most frequent serial reasons being early exits (\Code{break}/\Code{return}) and scalar recurrences. The \fieldUnk{} \textsc{Unknown} loops (\fieldPctUnk{}) are the population on which the hybrid procedure operates; their most frequent reasons are unresolved dependence directions and calls with unknown side effects (often STL member functions). The whole corpus is analysed in \fieldSeconds{}\,s, of which the pass itself takes \plPassSeconds{}\,s; the rest is IR generation.

% generated by hf/paper_extra.py - do not edit by hand

\begin{table}[htbp]
\caption{Loops of the 200-file application corpus (37,719 lines, 10 domains) through the deterministic phases.}
\label{tab:pipeline}
\centering
\small
\begin{tabular}{lrr}
\toprule
Pipeline phase & Loops & Share \\
\midrule
Phase 1: natural loops found (incl.\ system/library headers) & 2,187 &  \\
Phase 1: loops in system/library headers removed & $-$343 &  \\
Phase 1: project loops & 1,844 & 100.0\% \\
\midrule
Phase 2: proven \textsc{Parallel} (directive emitted, no model call) & 303 & 16.4\% \\
Phase 2: proven \textsc{Serial} (veto) & 592 & 32.1\% \\
Phase 2: \textsc{Unknown} $\rightarrow$ candidates for Phases 3--4 & \textbf{949} & \textbf{51.5\%} \\
\bottomrule
\end{tabular}
\end{table}


\subsection{RQ1: The LLVM-only baseline under an execution-validated ground truth}
\label{sec:rq-static}
Table~\ref{tab:fullset} compares the three sound compiler tools on all \nBenchLoops{} loops of HF-LoopBench.

% generated by hf/report.py from results/metrics.json - do not edit by hand
\begin{table}[htbp]
\caption{Compiler-based tools on all 990 target loops (95\% cluster-bootstrap CI).}
\label{tab:fullset}
\centering
\small
\resizebox{\linewidth}{!}{%
\begin{tabular}{lrrrrrrrr}
\toprule
Tool & TP & FP & TN & FN & Prec. & Rec. & F1 [95\% CI] & MCC \\
\midrule
HF-Static (LLVM 21) & 377 & 2 & 334 & 277 & 99.5 & 57.6 & 73.0 [67.7, 77.9] & 0.556 \\
GCC autopar & 233 & 0 & 336 & 421 & 100.0 & 35.6 & 52.5 [44.9, 59.9] & 0.398 \\
Polly & 66 & 7 & 329 & 588 & 90.4 & 10.1 & 18.2 [9.0, 27.1] & 0.145 \\
\bottomrule
\end{tabular}
}
\end{table}


HF-Static decides \nStaticSertoSer{} of the \nBenchSer{} serial loops correctly as \textsc{Serial} and accepts \fsStaticTP{} of the \nBenchPar{} parallel loops (precision \fsStaticPrec{}, recall \fsStaticRec{}, F1 \fsStaticFone{}). GCC is equally conservative but finds fewer loops (recall \fsGccRec{}), and Polly parallelizes only \fsPollyTP{} target loops because most kernels are not static control parts. The \nStaticFalsePar{} serial loops that HF-Static labels \textsc{Parallel} are the \Code{restrict} variants of a kernel that stages every element through a caller-provided scratch cell (\Code{*scratch = a[i]*2; b[i] = *scratch + 1;}): with \Code{restrict}, LLVM promotes the scratch cell to a register before the analysis runs, so the IR loop is genuinely independent, but no OpenMP clause can privatize \Code{*scratch} at source level. This is a limitation of IR-level analysis for source-level directives rather than of dependence testing. Polly's \fsPollyFP{} ``false positives'' fall outside our directive-only definition: six are variants of a column-wise integral-image nest that Polly parallelizes after transforming the nest, and one is an overlapping-argument variant, for which Polly emits code versioned on run-time alias checks.

The price of soundness is the \textsc{Unknown} class: \nStaticPartoUnk{} parallel and \nStaticSertoUnk{} serial loops are left undecided. The recorded reasons (Table~\ref{tab:unknown}) are dominated by pointer parameters that alias analysis cannot separate and by flattened multi-dimensional subscripts (\Code{a[i*m+j]}) for which DependenceAnalysis cannot resolve the direction at the outer level. The pass needs \staticUsPerLoop{}\,\textmu s per loop on average. In short, the LLVM-only configuration is safe but finds little more than half of the parallelism; the 96.6\% F1 of our preliminary version was an artefact of using the pass as its own oracle.

\begin{table}[htbp]
\caption{Reasons recorded by HF-Static for \textsc{Unknown} target loops of HF-LoopBench (a loop can have several).}
\label{tab:unknown}
\centering\small
% generated by hf/report.py
\begin{tabular}{lr}
\toprule
Reason reported by HF-Static & Loops \\
\midrule
dependence test inconclusive (may-alias or non-affine subscript) & 213 \\
possible loop-carried output dependence (direction not resolved) & 206 \\
possible loop-carried anti dependence (direction not resolved) & 159 \\
conditionally updated scalar (possible min/max/last-value reduction) & 75 \\
possible loop-carried flow dependence (direction not resolved) & 24 \\
\bottomrule
\end{tabular}

\end{table}

\subsection{RQ2: LLMs alone and the hallucination problem}
\label{sec:rq-ai}
Table~\ref{tab:evalset} lists all methods on the \nEvalLoops{}-loop evaluation set. Contrary to our preliminary version, current open-weight models are \emph{strong} classifiers when used alone: with prompt P2, GPT-OSS-120B reaches F1 \evAIOneTwentyFone{} and GPT-OSS-20B \evAITwentyFone{}, well above HF-Static (\evStaticFone{}) on the same loops, because they recognize stencils, gathers, privatizable temporaries and pointer-parameter DOALL loops that the compiler cannot prove.

% generated by hf/report.py from results/metrics.json - do not edit by hand
\begin{table}[htbp]
\caption{All methods on the LLM evaluation set (396 loops).}
\label{tab:evalset}
\centering
\small
\resizebox{\linewidth}{!}{%
\begin{tabular}{lrrrrrrrr}
\toprule
Method & TP & FP & TN & FN & Prec. & Rec. & F1 [95\% CI] & MCC \\
\midrule
HF-Static (LLVM 21) & 115 & 0 & 154 & 127 & 100.0 & 47.5 & 64.4 [58.8, 69.6] & 0.510 \\
GCC autopar & 73 & 0 & 154 & 169 & 100.0 & 30.2 & 46.3 [39.2, 53.2] & 0.379 \\
Polly & 24 & 3 & 151 & 218 & 88.9 & 9.9 & 17.8 [9.0, 26.3] & 0.154 \\
AI-only GPT-OSS-120B & 219 & 38 & 116 & 23 & 85.2 & 90.5 & 87.8 [84.6, 90.6] & 0.672 \\
Hybrid GPT-OSS-120B, no facts ($\tau^*$) & 215 & 2 & 152 & 27 & 99.1 & 88.8 & 93.7 [90.5, 96.4] & 0.858 \\
Hybrid GPT-OSS-120B, no call-site guard ($\tau^*$) & 136 & 7 & 147 & 106 & 95.1 & 56.2 & 70.6 [65.1, 75.9] & 0.524 \\
Hybrid GPT-OSS-120B ($\tau{=}0.5$) & 217 & 2 & 152 & 25 & 99.1 & 89.7 & 94.1 [91.1, 96.7] & 0.866 \\
Hybrid GPT-OSS-120B ($\tau^*$) & 136 & 0 & 154 & 106 & 100.0 & 56.2 & 72.0 [66.3, 77.1] & 0.577 \\
AI-only GPT-OSS-20B & 202 & 43 & 111 & 40 & 82.4 & 83.5 & 83.0 [79.2, 86.4] & 0.558 \\
Hybrid GPT-OSS-20B, no facts ($\tau^*$) & 212 & 5 & 149 & 30 & 97.7 & 87.6 & 92.4 [88.7, 95.4] & 0.826 \\
Hybrid GPT-OSS-20B, no call-site guard ($\tau^*$) & 209 & 34 & 120 & 33 & 86.0 & 86.4 & 86.2 [82.6, 89.4] & 0.644 \\
Hybrid GPT-OSS-20B ($\tau{=}0.5$) & 209 & 0 & 154 & 33 & 100.0 & 86.4 & 92.7 [89.1, 95.8] & 0.843 \\
Hybrid GPT-OSS-20B ($\tau^*$) & 209 & 0 & 154 & 33 & 100.0 & 86.4 & 92.7 [89.1, 95.8] & 0.843 \\
\bottomrule
\end{tabular}
}
\end{table}


They are, however, not safe. GPT-OSS-120B declares \evAIOneTwentyFP{} serial loops parallel and GPT-OSS-20B \evAITwentyFP{} (precision \evAIOneTwentyPrec{} and \evAITwentyPrec{}); each of these would be a data race. For the primary model, \nAIFPQ{} of the unsafe answers concern the overlapping-call variants: the function body is a textbook DOALL loop, the call site printed in the prompt (\Code{f(buf + 1, buf, n)}) makes it serial, and the model's explanation (``each iteration writes distinct \Code{y[i]}'') shows that the call site was ignored. The remaining \nAIFPnonQ{} are loop-carried flow dependences and scalar recurrences that were misread, a conflicting scatter and the shared scratch cell.

Table~\ref{tab:aidist} breaks the model's answers down by self-reported confidence. Most answers, correct or not, carry a confidence of at least 0.95, and every one of the \aiUnsafeTwenty{} unsafe answers of GPT-OSS-20B (\aiHiConfUnsafeOneTwenty{} of \aiUnsafeOneTwenty{} for GPT-OSS-120B) was given with a confidence of at least 0.9. Confidence therefore separates correct from incorrect answers only weakly, which limits what a threshold alone can achieve (Section~\ref{sec:rq-tau}). On the other side, both models systematically reject histogram-style array reductions (\Code{hist[px[i]]++}) that \Code{reduction(+:hist[:256])} makes legal (Section~\ref{sec:case-hist}).

% generated by hf/paper_extra.py - do not edit by hand

\begin{table}[htbp]
\caption{Answers of the LLMs alone (prompt P2, evaluation set) by self-reported confidence $c$. Unsafe: ``parallel'' answers for loops whose GT is serial; missed: ``serial'' answers for GT-parallel loops.}
\label{tab:aidist}
\centering
\small
\resizebox{\linewidth}{!}{%
\begin{tabular}{lrrrrrrrr}
\toprule
 & \multicolumn{4}{c}{GPT-OSS-20B} & \multicolumn{4}{c}{GPT-OSS-120B} \\ Confidence & ``parallel'' & unsafe & ``serial'' & missed & ``parallel'' & unsafe & ``serial'' & missed \\
\midrule
$c \ge 0.95$ & 245 & 43 & 145 & 37 & 251 & 38 & 136 & 22 \\
$0.90 \le c < 0.95$ & 0 & 0 & 5 & 2 & 6 & 0 & 3 & 1 \\
$0.80 \le c < 0.90$ & 0 & 0 & 1 & 1 & 0 & 0 & 0 & 0 \\
$c < 0.80$ & 0 & 0 & 0 & 0 & 0 & 0 & 0 & 0 \\
no parsable answer & 0 &  &  &  & 0 &  &  &  \\
\bottomrule
\end{tabular}
}
\end{table}


\subsection{RQ3: The hybrid procedure}
\label{sec:rq-hybrid}

\paragraph{Static veto} In the hybrid pipeline the model is asked only about the \nEvalUnknown{} \textsc{Unknown} loops of the evaluation set; every other decision comes from HF-Static. This removes all unsafe answers on loops with a proven dependence: for the primary model the veto overrides \nVetoSaves{} unsafe answers directly. It cannot remove the overlapping-call variants, for which HF-Static analyses the callee in isolation and reports \textsc{Unknown}: without a further guard, the hybrid with GPT-OSS-20B still emits \evHTwentyFP{} unsafe directives (row ``no call-site guard'' in Table~\ref{tab:evalset}), all of them on these loops.

\paragraph{Call-site guard} With the guard of Section~\ref{sec:static}, which flags exactly the \nCtxFlaggedAll{} overlapping-call instances of HF-LoopBench and no other kernel, the complete hybrid with GPT-OSS-20B ($\tau^*=\tauStarTwenty$) accepts \evHCTwentyTP{} of the \nEvalPar{} parallel loops of the evaluation set without a single unsafe directive (precision \evHCTwentyPrec{}, recall \evHCTwentyRec{}, F1 \evHCTwentyFone{} [\evHCTwentyFoneLo{}, \evHCTwentyFoneHi{}], MCC \evHCTwentyMCC{}); on the held-out 70\% test split the F1 is \evHCTwentyTestFone{} with \evHCTwentyTestFP{} false positives. Compared with the LLVM-only baseline, recall rises from \evStaticRec{} to \evHCTwentyRec{} at unchanged precision; compared with the LLM alone, all \evAITwentyFP{} unsafe answers disappear and F1 rises from \evAITwentyFone{} to \evHCTwentyFone{}.

\paragraph{Effect of the compiler facts} The row ``no facts'' answers the \textsc{Unknown} loops with the P2 answers (no compiler facts) under the same veto and guard. For GPT-OSS-20B the facts remove the \evHBTwentyFP{} unsafe answers that remain without them at a similar recall (\evHCTwentyRec{} vs.\ \evHBTwentyRec{}). For GPT-OSS-120B they do not help: the facts lower the confidence the model reports, so that at its calibrated threshold the complete system accepts fewer loops (F1 \evHCOneTwentyFone{}) than the no-facts variant (F1 \evHBOneTwentyFone{}, \evHBOneTwentyFP{} unsafe directives). Injecting compiler facts is therefore not uniformly beneficial; what makes the pipeline safe is the veto and the guard, and the threshold must be calibrated per model and prompt.

\subsubsection{Statistical significance}
Table~\ref{tab:mcnemar} reports exact McNemar tests on per-loop correctness. The complete hybrid is significantly more often correct than both single-modality baselines for both models ($p<0.001$); for GPT-OSS-20B there is no loop on which the LLVM-only baseline is correct and the hybrid is not.

% generated by hf/paper_extra.py - do not edit by hand

\begin{table}[htbp]
\caption{Exact McNemar tests on per-loop correctness (evaluation set, 396 loops). $b$: loops only the first method gets right; $c$: loops only the second gets right.}
\label{tab:mcnemar}
\centering
\footnotesize
\begin{tabular}{lrrr}
\toprule
Comparison & $b$ & $c$ & $p$ \\
\midrule
Hybrid GPT-OSS-120B vs.\ AI-only GPT-OSS-120B & 40 & 85 & $<0.001$ \\
Hybrid GPT-OSS-120B vs.\ LLVM-only & 21 & 0 & $<0.001$ \\
Hybrid GPT-OSS-120B vs.\ Hybrid GPT-OSS-120B w/o guard & 7 & 0 & 0.016 \\
AI-only GPT-OSS-120B vs.\ LLVM-only & 106 & 40 & $<0.001$ \\
Hybrid GPT-OSS-120B vs.\ Hybrid GPT-OSS-120B w/o facts & 2 & 79 & $<0.001$ \\
Hybrid GPT-OSS-120B vs.\ GCC & 67 & 4 & $<0.001$ \\
Hybrid GPT-OSS-20B vs.\ AI-only GPT-OSS-20B & 57 & 7 & $<0.001$ \\
Hybrid GPT-OSS-20B vs.\ LLVM-only & 94 & 0 & $<0.001$ \\
Hybrid GPT-OSS-20B vs.\ Hybrid GPT-OSS-20B w/o guard & 34 & 0 & $<0.001$ \\
AI-only GPT-OSS-20B vs.\ LLVM-only & 97 & 53 & $<0.001$ \\
Hybrid GPT-OSS-20B vs.\ Hybrid GPT-OSS-20B w/o facts & 9 & 7 & 0.804 \\
Hybrid GPT-OSS-20B vs.\ GCC & 136 & 0 & $<0.001$ \\
\bottomrule
\end{tabular}
\end{table}


\subsubsection{Confidence-threshold sensitivity (\texorpdfstring{$\tau$}{tau})}
\label{sec:rq-tau}
The threshold lets a user govern the model's influence (Tables~\ref{tab:sweep-gpt-oss-20b} and~\ref{tab:sweep-gpt-oss-120b}). For GPT-OSS-20B the complete system is already safe at $\tau=0.5$ on the development split, so calibration selects $\tau^*=\tauStarTwenty$; raising $\tau$ only costs recall (F1 \swHCFoneNinetyFiveTwenty{} at 0.95). For GPT-OSS-120B the development split contains one serial loop---an in-place running maximum, \Code{img[i] = max(img[i-1], img[i])}---that the model calls parallel with confidence above 0.9, so calibration selects $\tau^*=\tauStarOneTwenty$: the system becomes safe (\evHCOneTwentyFP{} unsafe directives) but its F1 drops to \evHCOneTwentyFone{}, whereas $\tau=0.5$ would give F1 \evHFiftyOneTwentyFone{} with \evHFiftyOneTwentyFP{} unsafe directives. As in our preliminary version, a high threshold buys safety; unlike in that version, the value is chosen on held-out kernels and its cost in recall is reported. The larger model is the better stand-alone classifier but the worse partner in the hybrid: its errors are confident, and confidence is the only signal the gate can use.

% generated by hf/report.py from results/metrics.json - do not edit by hand
\begin{table}[htbp]
\caption{Sensitivity to the confidence threshold $\tau$ (openai/gpt-oss-20b).}
\label{tab:sweep-gpt-oss-20b}
\centering
\small
\resizebox{\linewidth}{!}{%
\begin{tabular}{rrrrrrrrrrr}
\toprule
$\tau$ & AI-only FP & AI-only F1 & No facts FP & No facts F1 & No guard FP & No guard F1 & Full TP & Full FP & Full F1 & Full FP (dev) \\
\midrule
0.50 & 43 & 83.0 & 5 & 92.4 & 34 & 86.2 & 209 & 0 & 92.7 & 0 \\
0.60 & 43 & 83.0 & 5 & 92.4 & 34 & 86.2 & 209 & 0 & 92.7 & 0 \\
0.70 & 43 & 83.0 & 5 & 92.4 & 34 & 86.2 & 209 & 0 & 92.7 & 0 \\
0.80 & 43 & 83.0 & 5 & 92.4 & 34 & 86.2 & 209 & 0 & 92.7 & 0 \\
0.85 & 43 & 83.0 & 5 & 92.4 & 34 & 86.0 & 208 & 0 & 92.4 & 0 \\
0.90 & 43 & 83.0 & 5 & 92.4 & 31 & 84.6 & 200 & 0 & 90.5 & 0 \\
0.95 & 43 & 83.0 & 5 & 92.4 & 16 & 77.1 & 162 & 0 & 80.2 & 0 \\
0.99 & 33 & 68.7 & 3 & 85.2 & 4 & 66.3 & 122 & 0 & 67.0 & 0 \\
\bottomrule
\end{tabular}
}
\par\smallskip\footnotesize AI-only: P2 answers gated by $\tau$. No facts: veto + call-site guard + P2 answers. No guard: veto + H answers. Full: veto + call-site guard + H answers (complete system). F1 in \%.
\end{table}

% generated by hf/report.py from results/metrics.json - do not edit by hand
\begin{table}[htbp]
\caption{Sensitivity to the confidence threshold $\tau$ (openai/gpt-oss-120b).}
\label{tab:sweep-gpt-oss-120b}
\centering
\small
\resizebox{\linewidth}{!}{%
\begin{tabular}{rrrrrrrrrrr}
\toprule
$\tau$ & AI-only FP & AI-only F1 & No facts FP & No facts F1 & No guard FP & No guard F1 & Full TP & Full FP & Full F1 & Full FP (dev) \\
\midrule
0.50 & 38 & 87.8 & 2 & 95.1 & 37 & 87.5 & 217 & 2 & 94.1 & 1 \\
0.60 & 38 & 87.8 & 2 & 95.1 & 37 & 87.5 & 217 & 2 & 94.1 & 1 \\
0.70 & 38 & 87.8 & 2 & 95.1 & 37 & 87.0 & 215 & 2 & 93.7 & 1 \\
0.80 & 38 & 87.8 & 2 & 95.1 & 37 & 85.9 & 210 & 2 & 92.5 & 1 \\
0.85 & 38 & 87.8 & 2 & 95.1 & 35 & 84.8 & 204 & 1 & 91.3 & 1 \\
0.90 & 38 & 87.8 & 2 & 95.1 & 28 & 82.9 & 191 & 1 & 88.0 & 1 \\
0.95 & 38 & 86.4 & 2 & 93.7 & 7 & 70.6 & 136 & 0 & 72.0 & 0 \\
0.99 & 30 & 74.9 & 0 & 87.7 & 0 & 64.4 & 115 & 0 & 64.4 & 0 \\
\bottomrule
\end{tabular}
}
\par\smallskip\footnotesize AI-only: P2 answers gated by $\tau$. No facts: veto + call-site guard + P2 answers. No guard: veto + H answers. Full: veto + call-site guard + H answers (complete system). F1 in \%.
\end{table}


\subsubsection{All loops of HF-LoopBench}
\label{sec:rq-full}
The evaluation set contains one element type and no \Code{restrict} variant per kernel. To score the complete pipeline on the whole benchmark, the remaining \textsc{Unknown} loops were answered later with prompt H at the frozen thresholds. Table~\ref{tab:fullhybrid} compares the result with the compiler tools on all \nBenchLoops{} loops. With GPT-OSS-20B the hybrid accepts \fhTwentyTP{} of the \nBenchPar{} parallel loops (recall \fhTwentyRec{}, F1 \fhTwentyFone{} [\fhTwentyFoneLo{}, \fhTwentyFoneHi{}]) with \fhTwentyFP{} unsafe directives\fhTwentyFPNote{} Two of them are the \Code{restrict} scratch-cell loops that HF-Static itself accepts (Section~\ref{sec:rq-static}); the other two are \Code{float} variants of a Bellman--Ford relaxation and of a time-stepping loop for which the model answered ``parallel'' with confidence 0.9, although it had answered ``serial'' with confidence 0.85--0.99 for the textually equivalent \Code{double} variants of the same kernels. The model is thus not even consistent across element types, which no threshold calibrated on other loops can anticipate. On the \fhOutN{} loops outside the evaluation set, which played no role in calibration, the hybrid's F1 is \fhTwentyOutFone{} against \fhStaticOutFone{} for HF-Static. With GPT-OSS-120B at its conservative $\tau^*$ the hybrid reaches F1 \fhOneTwentyFone{} with \fhOneTwentyFP{} unsafe directives (both the scratch-cell loops), whereas $\tau=0.5$ would give \fhOneTwentyFiftyFone{} with \fhOneTwentyFiftyFP{}. The model was queried for \fhTwentyQueried{} of the \nBenchLoops{} loops (\fhQueriedPct{}).

% generated by hf/paper_extra.py - do not edit by hand

\begin{table}[htbp]
\caption{All 990 loops of HF-LoopBench: sound compiler tools versus the complete hybrid (thresholds frozen at the values calibrated on the evaluation set's development split).}
\label{tab:fullhybrid}
\centering
\small
\resizebox{\linewidth}{!}{%
\begin{tabular}{lrrrrrrrr}
\toprule
Method & TP & FP & TN & FN & Prec. & Rec. & F1 [95\% CI] & MCC \\
\midrule
LLVM-only (HF-Static) & 377 & 2 & 334 & 277 & 99.5 & 57.6 & 73.0 [67.7, 77.9] & 0.556 \\
GCC autopar & 233 & 0 & 336 & 421 & 100.0 & 35.6 & 52.5 [44.9, 59.9] & 0.398 \\
Polly & 66 & 7 & 329 & 588 & 90.4 & 10.1 & 18.2 [9.0, 27.1] & 0.145 \\
Hybrid GPT-OSS-120B ($\tau{=}0.5$) & 577 & 8 & 328 & 77 & 98.6 & 88.2 & 93.1 & 0.827 \\
\textbf{Hybrid GPT-OSS-120B} ($\tau^*{=}0.95$) & 411 & 2 & 334 & 243 & 99.5 & 62.8 & 77.0 [72.2, 81.7] & 0.598 \\
Hybrid GPT-OSS-20B ($\tau{=}0.5$) & 566 & 4 & 332 & 88 & 99.3 & 86.5 & 92.5 & 0.818 \\
\textbf{Hybrid GPT-OSS-20B} ($\tau^*{=}0.5$) & 566 & 4 & 332 & 88 & 99.3 & 86.5 & 92.5 [89.0, 95.5] & 0.818 \\
\bottomrule
\end{tabular}
}
\end{table}


\subsubsection{Loop-structure extension (HF-LoopBench-X)}
\label{sec:rq-ext}
Table~\ref{tab:ext} reports all methods on the \xLoops{} loops of HF-LoopBench-X, with every hybrid threshold frozen at the value calibrated on HF-LoopBench, and Table~\ref{tab:ext-classes} breaks the results down by loop structure. \extText{}

The two unsafe directives of the GPT-OSS-20B hybrid are the global and pointer variants of \Code{for (i = 0; i < n - k; i++) a[i] = a[i + k] * c;}, an anti dependence whose distance \Code{k} is a run-time parameter (3 at the only call site). HF-Static cannot resolve the direction of a symbolic distance and reports \textsc{Unknown}; the model answered ``parallel'' with confidence 0.9 and 0.95 (``each iteration writes distinct \Code{a[i]} and reads \Code{a[i+k]}''), while it answered ``serial'' for the \Code{restrict} variant. This is a residual risk that the veto does not cover: the fact that decides legality (the sign and size of \Code{k}) is known only at run time. A run-time version check (\Code{k == 0 || k >= n}) or a veto on unresolved dependences with symbolic distance would remove it at the cost of some recall; we did not add either after seeing the results, so as not to tune on the test set. The compiler facts also made GPT-OSS-20B reject the parallel inner loop of the vertical accumulation (``loop-carried dependence on previous \Code{k-1}''), confusing the level that carries the dependence with the target level. Apart from the disjoint-halves calls that the guard vetoes, the inner-level targets are where the LLM alone is clearly better than the hybrid (Table~\ref{tab:ext-classes}). With Qwen3.8-27B, the frozen hybrid makes no unsafe decision on the extension (F1 \xHQwenFone{}).

% generated by hf/ext_analysis.py from results/ext - do not edit by hand

\begin{table}[htbp]
\caption{HF-LoopBench-X (121 loops, written after the main study; hybrid thresholds frozen at the values calibrated on the main benchmark).}
\label{tab:ext}
\centering\small
\resizebox{\linewidth}{!}{%
\begin{tabular}{lrrrrrrrr}
\toprule
Method & TP & FP & TN & FN & Prec. & Rec. & F1 [95\% CI] & MCC \\
\midrule
HF-Static (LLVM 21) & 45 & 0 & 41 & 35 & 100.0 & 56.2 & 72.0 [58.8, 82.7] & 0.551 \\
GCC autopar & 32 & 0 & 41 & 48 & 100.0 & 40.0 & 57.1 [40.7, 71.2] & 0.429 \\
Polly & 9 & 3 & 38 & 71 & 75.0 & 11.2 & 19.6 [0.0, 39.5] & 0.062 \\
AI-only GPT-OSS-20B & 77 & 11 & 30 & 3 & 87.5 & 96.2 & 91.7 [87.4, 95.1] & 0.738 \\
Hybrid GPT-OSS-20B, no call-site guard ($\tau^*{=}0.5$) & 73 & 9 & 32 & 7 & 89.0 & 91.2 & 90.1 [83.0, 95.3] & 0.702 \\
Hybrid GPT-OSS-20B ($\tau{=}0.5$) & 65 & 2 & 39 & 15 & 97.0 & 81.2 & 88.4 [80.3, 94.2] & 0.727 \\
Hybrid GPT-OSS-20B ($\tau^*{=}0.5$) & 65 & 2 & 39 & 15 & 97.0 & 81.2 & 88.4 [80.3, 94.2] & 0.727 \\
AI-only GPT-OSS-120B & 77 & 9 & 32 & 3 & 89.5 & 96.2 & 92.8 [89.3, 96.0] & 0.776 \\
Hybrid GPT-OSS-120B, no call-site guard ($\tau^*{=}0.95$) & 57 & 2 & 39 & 23 & 96.6 & 71.2 & 82.0 [71.6, 90.0] & 0.628 \\
Hybrid GPT-OSS-120B ($\tau{=}0.5$) & 64 & 6 & 35 & 16 & 91.4 & 80.0 & 85.3 [77.5, 91.4] & 0.627 \\
Hybrid GPT-OSS-120B ($\tau^*{=}0.95$) & 51 & 1 & 40 & 29 & 98.1 & 63.7 & 77.3 [66.0, 86.1] & 0.586 \\
AI-only Qwen3.8-27B & 72 & 7 & 34 & 8 & 91.1 & 90.0 & 90.6 [84.6, 95.2] & 0.725 \\
Hybrid Qwen3.8-27B, no call-site guard ($\tau^*{=}0.5$) & 75 & 6 & 35 & 5 & 92.6 & 93.8 & 93.2 [87.5, 97.4] & 0.796 \\
Hybrid Qwen3.8-27B ($\tau{=}0.5$) & 67 & 0 & 41 & 13 & 100.0 & 83.8 & 91.2 [85.0, 96.0] & 0.797 \\
Hybrid Qwen3.8-27B ($\tau^*{=}0.5$) & 67 & 0 & 41 & 13 & 100.0 & 83.8 & 91.2 [85.0, 96.0] & 0.797 \\
\bottomrule
\end{tabular}
}
\end{table}

% generated by hf/ext_analysis.py from results/ext - do not edit by hand

\begin{table}[htbp]
\caption{HF-LoopBench-X: accuracy (\%) per loop-structure class. Unk = loops HF-Static leaves \textsc{Unknown}; AI = LLM alone (P2); Hyb = complete hybrid at the frozen $\tau^*$.}
\label{tab:ext-classes}
\centering\scriptsize
\resizebox{\linewidth}{!}{%
\begin{tabular}{llrrrrrrrrrrr}
\toprule
Class & GT & $n$ & Unk & HF-S & GCC & Polly & AI 20B & AI 120B & AI Q27B & Hyb 20B & Hyb 120B & Hyb Q27B \\
\midrule
Called with disjoint halves & P & 8 & 8 & 0 & 25 & 0 & 100 & 88 & 100 & 0 & 0 & 0 \\
Decrementing loop & P & 3 & 1 & 67 & 100 & 0 & 100 & 100 & 100 & 100 & 67 & 100 \\
Disjoint in-place / modulo writes & P & 9 & 3 & 67 & 44 & 0 & 100 & 100 & 100 & 100 & 67 & 100 \\
Inner loop of a nest & P & 9 & 7 & 22 & 22 & 33 & 100 & 78 & 67 & 56 & 44 & 78 \\
Non-unit step & P & 3 & 1 & 67 & 67 & 0 & 100 & 100 & 100 & 100 & 67 & 100 \\
Perfect nest, collapse(2) & P & 3 & 0 & 100 & 0 & 100 & 100 & 100 & 100 & 100 & 100 & 100 \\
Permutation scatter & P & 3 & 3 & 0 & 0 & 0 & 33 & 100 & 0 & 0 & 0 & 0 \\
Privatisable temporary array & P & 3 & 1 & 67 & 67 & 0 & 100 & 100 & 100 & 100 & 100 & 100 \\
Product / nested reduction & P & 9 & 0 & 100 & 0 & 0 & 89 & 100 & 89 & 100 & 100 & 100 \\
Same-iteration read-after-write & P & 3 & 1 & 67 & 67 & 0 & 100 & 100 & 100 & 100 & 67 & 100 \\
Strided access & P & 6 & 1 & 83 & 83 & 0 & 100 & 100 & 100 & 100 & 83 & 100 \\
Triangular nest & P & 6 & 4 & 33 & 33 & 50 & 100 & 100 & 100 & 100 & 33 & 100 \\
Unsigned induction variable & P & 3 & 1 & 67 & 67 & 0 & 100 & 100 & 100 & 100 & 100 & 100 \\
Walking pointers (linear) & P & 3 & 1 & 67 & 67 & 0 & 100 & 100 & 67 & 100 & 100 & 100 \\
if/else, mask, continue & P & 9 & 3 & 67 & 44 & 0 & 100 & 100 & 100 & 100 & 78 & 100 \\
Anti dependence, run-time offset & S & 3 & 3 & 100 & 100 & 100 & 33 & 67 & 100 & 33 & 67 & 100 \\
Arg-min (no UDR) & S & 3 & 3 & 100 & 100 & 100 & 100 & 100 & 100 & 100 & 100 & 100 \\
Called with overlapping arguments & S & 8 & 8 & 100 & 100 & 100 & 0 & 0 & 12 & 100 & 100 & 100 \\
Carried flow (distance 1/2, cross-stmt) & S & 9 & 0 & 100 & 100 & 100 & 89 & 100 & 100 & 100 & 100 & 100 \\
Cond. / ring-buffer / list recurrence & S & 9 & 6 & 100 & 100 & 100 & 100 & 100 & 100 & 100 & 100 & 100 \\
Outer loop carrying the dependence & S & 6 & 6 & 100 & 100 & 50 & 100 & 100 & 100 & 100 & 100 & 100 \\
Wavefront (inner loop) & S & 3 & 1 & 100 & 100 & 100 & 100 & 100 & 100 & 100 & 100 & 100 \\
\bottomrule
\end{tabular}
}
\end{table}


\paragraph{A soundness bug found by the extension} Before the fix described in Section~\ref{sec:static}, HF-Static labelled the global and \Code{restrict} variants of the linked-list walk (\Code{out[k] = val[p]; p = nxt[p];}) \textsc{Parallel}. Its rule for header $\phi$-nodes treated a value forwarded from a load of read-only memory as benign---correct for the redundant loads that GVN and LICM leave in a loop, wrong when the load address depends on the carried value itself. The execution oracle confirmed the race. We added the address condition, rebuilt the pass and re-ran the static stage on all benchmarks and on the application corpus: \fixText{}

\subsubsection{Categorization by dependence class}
Table~\ref{tab:classes} breaks the HF-LoopBench evaluation set down by dependence class for the primary model. The hybrid matches the LLM on stencils, gathers and privatizable scalars, where the compiler is weak, and matches the compiler on every serial class. Two classes remain open: array reductions, which no method accepts, and disjoint non-affine writes, on which the compiler facts make the model more conservative than it is without them.

% generated by hf/report.py from results/metrics.json - do not edit by hand
\begin{table}[htbp]
\caption{Accuracy (\%) per dependence class on the evaluation set (primary model).}
\label{tab:classes}
\centering
\small
\resizebox{\linewidth}{!}{%
\begin{tabular}{lcrrrrrrr}
\toprule
Dependence class & GT & Loops & Static UNK & Static & GCC & AI-only & Hybrid & Hybrid+ctx \\
\midrule
DOALL (affine) & P & 118 & 56 & 52.5 & 44.1 & 98.3 & 98.3 & 98.3 \\
Scalar reduction & P & 46 & 12 & 73.9 & 8.7 & 73.9 & 97.8 & 97.8 \\
Array reduction & P & 22 & 22 & 0.0 & 0.0 & 4.5 & 0.0 & 0.0 \\
Stencil & P & 22 & 16 & 27.3 & 22.7 & 95.5 & 95.5 & 95.5 \\
Disjoint non-affine writes & P & 14 & 10 & 28.6 & 14.3 & 78.6 & 64.3 & 64.3 \\
Gather / LUT & P & 10 & 5 & 50.0 & 60.0 & 100.0 & 100.0 & 100.0 \\
Privatisable scalar & P & 6 & 3 & 50.0 & 50.0 & 100.0 & 100.0 & 100.0 \\
Linear induction & P & 2 & 1 & 50.0 & 50.0 & 100.0 & 100.0 & 100.0 \\
Permutation scatter & P & 2 & 2 & 0.0 & 0.0 & 50.0 & 0.0 & 0.0 \\
Call-site aliasing & S & 35 & 35 & 100.0 & 100.0 & 5.7 & 2.9 & 100.0 \\
Carried flow/anti & S & 32 & 14 & 100.0 & 100.0 & 84.4 & 100.0 & 100.0 \\
Compaction / counter & S & 10 & 10 & 100.0 & 100.0 & 100.0 & 100.0 & 100.0 \\
Conflicting scatter & S & 12 & 12 & 100.0 & 100.0 & 91.7 & 100.0 & 100.0 \\
Early exit / non-canonical & S & 8 & 0 & 100.0 & 100.0 & 100.0 & 100.0 & 100.0 \\
Order-dependent call & S & 2 & 0 & 100.0 & 100.0 & 100.0 & 100.0 & 100.0 \\
Scalar recurrence & S & 50 & 4 & 100.0 & 100.0 & 94.0 & 100.0 & 100.0 \\
Shared scratch & S & 1 & 1 & 100.0 & 100.0 & 0.0 & 100.0 & 100.0 \\
Wavefront & S & 4 & 4 & 100.0 & 100.0 & 100.0 & 100.0 & 100.0 \\
\bottomrule
\end{tabular}
}
\end{table}


\subsection{RQ4: Caching, cost and prompts}

\subsubsection{Structural answer cache}
\label{sec:rq-cache}
Table~\ref{tab:cache} replays the recorded hybrid answers of the primary model through the structural cache in source order. Exact structural duplicates are rare among the queried loops (\cacheHitExact{} hit rate), because instances of one base kernel that reach the model differ in their data-access form. Lowering the similarity threshold to the value used in our preliminary version (0.85) raises the hit rate to \cacheHitEightyFive{} but turns \cacheFPCEightyFive{} correct \textsc{Serial} answers into unsafe \textsc{Parallel} answers inherited from a similar-looking loop; at 0.95 the cache avoids \cacheHitNinetyFive{} of the requests with \cacheFPCNinetyFive{} unsafe answers. Similarity of code text is not similarity of dependence structure---a one-token change from \Code{a[i-1]} to \Code{b[i-1]} flips the verdict---so the deployed configuration uses exact-match caching only. The 99\% ``warm-run'' hit rate of our preliminary version measured the re-analysis of unchanged files, which any content hash achieves; in a CI/CD setting this remains the main saving, since unchanged functions need not be re-queried at all.

% generated by hf/report.py from results/metrics.json - do not edit by hand
\begin{table}[htbp]
\caption{Structural cache replayed over the recorded hybrid answers.}
\label{tab:cache}
\centering
\small
\resizebox{\linewidth}{!}{%
\begin{tabular}{rrrrrrrrr}
\toprule
Sim. & Queries & Hits & Hit rate & Acc. fresh & Acc. cached & Changed & FP fresh & FP cached \\
\midrule
1.00 & 172 & 4 & 2.3 & 80.8 & 80.8 & 0 & 0 & 0 \\
0.95 & 172 & 11 & 6.4 & 80.8 & 80.2 & 1 & 0 & 0 \\
0.90 & 172 & 17 & 9.9 & 80.8 & 80.2 & 1 & 0 & 0 \\
0.85 & 172 & 35 & 20.3 & 80.8 & 77.9 & 9 & 0 & 4 \\
0.80 & 172 & 59 & 34.3 & 80.8 & 79.1 & 15 & 0 & 7 \\
\bottomrule
\end{tabular}
}
\end{table}


\subsubsection{Token cost and financial viability}
\label{sec:rq-cost}
Table~\ref{tab:costall} lists every model run of this study with its measured token usage and the cost at the provider's on-demand list prices (all runs actually used the free tier). Reasoning tokens are billed as output tokens; the column ``cached'' applies the provider's 50\% discount for prompt tokens served from its prompt cache, which matters because the system prompt is identical across batches. Because the hybrid queries the model only for the \projFracTwenty{} of loops that remain undecided after the static phase and the guard, it needs \projHTwenty{}\,USD per 1{,}000 analysed loops with GPT-OSS-20B instead of \projAITwenty{}\,USD for the LLM-only approach (\projSaveTwenty{} less, Figure~\ref{fig:cost}), although each hybrid prompt is longer because of the compiler facts (\tokTwentyHUnk{} vs.\ \tokTwentyPTwoEval{} tokens per loop). The corresponding figures for GPT-OSS-120B are \projHOneTwenty{} and \projAIOneTwenty{}\,USD. Under the free tier's daily token limit the same ratio determines throughput: \dayCapAITwenty{} loops per day for LLM-only classification versus \dayCapHTwenty{} analysed loops per day for the hybrid. For the 200-file corpus, LLM-only classification of all \fieldLoops{} loops would cost about \fieldAIUSD{}\,USD with GPT-OSS-20B, the hybrid about \fieldHUSD{}\,USD. In absolute terms all of these amounts are small; the entire experimental campaign of this article---\totRequests{} requests and \totTokens{} tokens---would have cost \totUSD{}\,USD at list prices. The dominant cost of the pipeline is therefore not money but latency and rate limits: a request of eight loops takes \latTwentyHUnk{}\,s with GPT-OSS-20B, whereas HF-Static needs \staticUsPerLoop{}\,\textmu s per loop.

% generated by hf/paper_extra.py - do not edit by hand

\begin{table}[htbp]
\caption{Every LLM run of this study: measured token usage and cost at Groq on-demand list prices (GPT-OSS-120B \$0.15/\$0.60, GPT-OSS-20B \$0.075/\$0.30, Qwen3.8-27B \$0.80/\$4.00 per million input/output tokens; ``cached'' applies the 50\% discount for cached prompt tokens of the GPT-OSS models). Output tokens include reasoning tokens (in parentheses).}
\label{tab:costall}
\centering
\scriptsize
\resizebox{\linewidth}{!}{%
\begin{tabular}{lllrrrrrrrr}
\toprule
Suite & Model & Prompt & Set & Loops & Req. & Input tok. & Output tok. (reas.) & Tok./loop & USD & USD cached \\
\midrule
Main & GPT-OSS-120B & H & \textsc{Unk} (eval.) & 207 & 26 & 65,062 & 15,257 (7,497) & 388 & 0.0189 & 0.0185 \\
Main & GPT-OSS-120B & H & \textsc{Unk} (rest) & 212 & 27 & 68,170 & 16,753 (8,946) & 401 & 0.0203 & 0.0200 \\
Main & GPT-OSS-120B & P2 & eval.\ set & 396 & 50 & 89,327 & 29,441 (15,004) & 300 & 0.0311 & 0.0305 \\
Main & GPT-OSS-20B & H & \textsc{Unk} (eval.) & 207 & 26 & 65,062 & 13,313 (5,435) & 379 & 0.0089 & 0.0088 \\
Main & GPT-OSS-20B & H & \textsc{Unk} (rest) & 212 & 27 & 68,170 & 13,709 (5,654) & 386 & 0.0092 & 0.0091 \\
Main & GPT-OSS-20B & P2 & eval.\ set & 396 & 50 & 89,327 & 23,019 (8,726) & 284 & 0.0136 & 0.0134 \\
Main & Qwen3.8-27B & P1 & ablation & 120 & 15 & 26,110 & 32,705 (0) & 490 & 0.1517 & 0.1517 \\
Main & Qwen3.8-27B & P2 & ablation & 120 & 15 & 29,485 & 5,621 (0) & 293 & 0.0461 & 0.0461 \\
Main & Qwen3.8-27B & P2 & ablation, run 2 & 120 & 15 & 29,485 & 5,604 (0) & 292 & 0.0460 & 0.0460 \\
Main & Qwen3.8-27B & P3 & ablation & 120 & 15 & 33,490 & 10,722 (0) & 368 & 0.0697 & 0.0697 \\
Main & Qwen3.8-27B & P4 & ablation & 120 & 15 & 33,805 & 5,986 (0) & 332 & 0.0510 & 0.0510 \\
X & GPT-OSS-20B & P2 & all & 121 & 16 & 25,635 & 7,547 (3,136) & 274 & 0.0042 & 0.0042 \\
X & GPT-OSS-20B & H & \textsc{Unk} & 63 & 8 & 18,695 & 3,846 (1,613) & 358 & 0.0026 & 0.0026 \\
X & GPT-OSS-120B & P2 & all & 121 & 16 & 25,635 & 10,759 (6,370) & 301 & 0.0103 & 0.0103 \\
X & GPT-OSS-120B & H & \textsc{Unk} & 63 & 8 & 18,695 & 5,063 (2,729) & 377 & 0.0058 & 0.0058 \\
X & Qwen3.8-27B & P2 & all & 121 & 16 & 27,405 & 5,844 (0) & 275 & 0.0453 & 0.0453 \\
X & Qwen3.8-27B & H & \textsc{Unk} & 63 & 8 & 20,066 & 3,122 (0) & 368 & 0.0285 & 0.0285 \\
\bottomrule
\end{tabular}
}
\end{table}

% generated by hf/paper_extra.py - do not edit by hand

\begin{figure}[htbp]
\centering
\begin{tikzpicture}
\begin{axis}[
    ybar, bar width=0.7cm, width=11cm, height=6.2cm, ymin=0, scaled y ticks=false, yticklabel style={/pgf/number format/fixed, /pgf/number format/precision=2},
    ylabel={USD per 1{,}000 analysed loops}, symbolic x coords={GPT-OSS-120B,GPT-OSS-20B}, xtick=data,
    nodes near coords={\pgfmathprintnumber[fixed,precision=3]\pgfplotspointmeta}, nodes near coords align={vertical},
    legend style={at={(0.5,-0.18)}, anchor=north, legend columns=1, font=\footnotesize}, enlarge x limits=0.45,
    ymajorgrids=true, grid style={dashed, gray!30}]
\addplot[fill=pipelineMagenta!70, draw=black] coordinates {(GPT-OSS-120B,0.0784) (GPT-OSS-20B,0.0344)};
\addplot[fill=pipelineBlue!70, draw=black] coordinates {(GPT-OSS-120B,0.0397) (GPT-OSS-20B,0.0186)};
\legend{LLM-only (every loop queried), \SystemName{} (only undecided loops queried)}
\end{axis}
\end{tikzpicture}
\caption{Measured token cost per 1{,}000 analysed loops at on-demand list prices (evaluation set).}
\label{fig:cost}
\end{figure}


\subsubsection{Prompt engineering and ablation}
\label{sec:rq-prompt}
Table~\ref{tab:ablation} compares the four prompt variants with Qwen3.8-27B on \nAblLoops{} loops, using real API calls, together with the P2 results of the GPT-OSS models on the same loops. \ablationText{}

% generated by hf/report.py from results/metrics.json - do not edit by hand
\begin{table}[htbp]
\caption{Prompt ablation (AI-only, no compiler facts).}
\label{tab:ablation}
\centering
\small
\resizebox{\linewidth}{!}{%
\begin{tabular}{lllrrrrrrr}
\toprule
Model & Prompt & Set & Loops & Parse fail & FP & Prec. & Rec. & F1 & Tok./loop \\
\midrule
qwen3.8-27b & P1 & ablation & 120 & 12.5 & 1 & 97.8 & 54.9 & 70.3 & 490 \\
qwen3.8-27b & P2 & ablation & 120 & 0.0 & 0 & 100.0 & 82.9 & 90.7 & 293 \\
qwen3.8-27b & P3 & ablation & 120 & 0.0 & 0 & 100.0 & 84.1 & 91.4 & 368 \\
qwen3.8-27b & P4 & ablation & 120 & 0.0 & 0 & 100.0 & 84.1 & 91.4 & 332 \\
gpt-oss-120b & P2 & ablation(eval subset) & 120 & -- & 1 & 98.7 & 90.2 & 94.3 & -- \\
gpt-oss-20b & P2 & ablation(eval subset) & 120 & -- & 4 & 94.6 & 85.4 & 89.7 & -- \\
\bottomrule
\end{tabular}
}
\end{table}


\paragraph{Run-to-run consistency} \consistencyText{}

\subsection{RQ5: Performance impact}
\label{sec:rq-perf}
We built every GT-parallel kernel of HF-LoopBench with 64$\times$ larger default problem sizes at \Code{-O2} and measured serial and OpenMP run time with \perfThreads{} threads (minimum of three runs). Of the \nPerfValid{} kernels with correct output, \perfFracGtOne{} run faster in parallel; the geometric-mean speedup is \perfGeo{}$\times$ (median \perfMedian{}$\times$, maximum \perfMax{}$\times$; Figure~\ref{fig:perf}). Many kernels touch only a few megabytes and finish in milliseconds, so thread start-up and memory bandwidth limit the gains; the figures characterize the benchmark rather than a tuned application. A practical comparison of the tools is the fraction of the total achievable time saving $\sum_L \max(0, T_\mathrm{ser}(L)-T_\mathrm{par}(L))$ that each tool's accepted loops capture: on all loops, HF-Static captures \covFullStatic{}, GCC \covFullGcc{}, Polly \covFullPolly{} and the complete hybrid with GPT-OSS-20B \covFullHTwenty{}. On HF-LoopBench-X the geometric-mean speedup of the \xPerfN{} parallel kernels is \xPerfGeo{}$\times$.

% generated by hf/paper_extra.py - do not edit by hand

\begin{figure}[htbp]
\centering
\begin{tikzpicture}
\begin{axis}[
    xbar, bar width=0.28cm, width=10cm, height=9cm, xmin=0,
    xlabel={Geometric-mean speedup, 8 threads}, symbolic y coords={bioinformatics,graph analytics,scientific misc,geometry graphics,networking,medical imaging,stencil pde,databases,cryptography,image processing,signal processing,linear algebra,machine learning,quantum simulation,finance,climate weather,molecular dynamics,sorting searching}, ytick=data, y tick label style={font=\footnotesize},
    nodes near coords={\pgfmathprintnumber[fixed,precision=2]\pgfplotspointmeta}, nodes near coords align={horizontal},
    xmajorgrids=true, grid style={dashed, gray!30}, enlarge y limits=0.04, extra x ticks={1}, extra x tick style={grid=major, grid style={black}}]
\addplot[fill=pipelineBlue!70, draw=black] coordinates {(3.903,sorting searching) (3.325,molecular dynamics) (3.055,climate weather) (2.707,finance) (2.678,quantum simulation) (2.228,machine learning) (2.127,linear algebra) (2.085,signal processing) (1.738,image processing) (1.704,cryptography) (1.689,databases) (1.533,stencil pde) (1.365,medical imaging) (1.286,networking) (1.280,geometry graphics) (1.102,scientific misc) (0.994,graph analytics) (0.815,bioinformatics)};
\end{axis}
\end{tikzpicture}
\caption{Speedup of the GT-parallel kernels of HF-LoopBench per domain (Apple M1 Pro, \Code{-O2}, minimum of three runs; kernels with wrong \Code{-O2} output excluded).}
\label{fig:perf}
\end{figure}

